\documentclass[a4paper, 12pt]{article}

\usepackage[french]{babel}
\usepackage[T1]{fontenc}
\usepackage[utf8]{inputenc}
\usepackage{amsmath, amssymb}
\usepackage{geometry}
\usepackage{array}
\usepackage{enumitem}
\usepackage{tikz}
\usepackage{pgfplots}
\pgfplotsset{compat=1.18}
\usetikzlibrary{angles, quotes, calc}

\geometry{top=2cm, bottom=2cm, left=2cm, right=2cm}

% CONFIGURATION
\newcommand{\sujet}{1}
\newcommand{\classe}{Première Spécialité}
\newcommand{\duree}{2 heures}
\newcommand{\datedevoir}{2026}
\newcommand{\theme}{Bac Blanc - Mathématiques}

% Cadre réponse Automatisme
\newcommand{\caseqcm}{
    \hfill \fbox{\textbf{Réponse :}\hspace{1cm}}
    \vspace{0.4cm}
}

% Cadre réponse simple
\newcommand{\reponse}{\par\vspace{0.3cm}
\framebox[\linewidth]{\begin{minipage}{\dimexpr\linewidth-2\fboxsep}
\vspace{2.5cm}\hspace{0.5cm}
\end{minipage}}\vspace{0.4cm}
}

% Cadre réponse double
\newcommand{\reponseDouble}{\par\vspace{0.3cm}
\framebox[\linewidth]{\begin{minipage}{\dimexpr\linewidth-2\fboxsep}
\vspace{5cm}\hspace{0.5cm}
\end{minipage}}\vspace{0.4cm}
}

\begin{document}

\begin{center}\Large\textbf{Bac blanc ( EAM 1ère  - Sujet \sujet)}\end{center}

\begin{tabular}{|p{8cm}|p{8cm}|}
\hline
\textbf{Nom et Prénom :} & \textbf{Date : \datedevoir} \\
\framebox[7.5cm]{\hspace{0.1cm}\begin{minipage}{7.3cm}\vspace{1.5cm}\end{minipage}} & \\
\hline
\end{tabular}

\begin{center}\textbf{\classe \ -- \ \theme}\end{center}

\textbf{Durée : \duree \ -- Calculatrice interdite en Partie 1}

\section*{Partie 1 : Automatismes (12 questions - 6 points)}

\begin{enumerate}[leftmargin=*, itemsep=0.4cm]
\item \textbf{Question :} Mettre $\frac{1}{3}-\frac{2-x}{2}$ sous la forme $\frac{a+bx}{c}$.
\begin{itemize}[label={}]
\item a) $\frac{-1+x}{1}$ \quad b) $\frac{2-3x}{6}$ \quad c) $\frac{3-x}{5}$ \quad d) $\frac{-4+3x}{6}$ \caseqcm
\end{itemize}

\item \textbf{Question :} Calculer $B = \frac{a}{c} + \frac{1/2}{b}$ pour $a = \frac{1}{6}, b = \frac{1}{2}, c = \frac{1}{3}$.
\begin{itemize}[label={}]
\item a) $\frac{1}{6}$ \quad b) $\frac{3}{2}$ \quad c) $\frac{1}{2}$ \quad d) 1 \caseqcm
\end{itemize}

\item \textbf{Question :} Soit la relation $C = \frac{2}{x} + \frac{3}{y}$. On peut alors affirmer que :
\begin{itemize}[label={}]
\item a) $x = 2(\frac{1}{C} - \frac{y}{3})$ \quad b) $x = \frac{2y}{Cy-3}$ \quad c) $x = \frac{5-Cy}{C}$ \caseqcm
\end{itemize}

\item \textbf{Question :} Proportion des 18-25 ans (3/4 mineurs, 1/3 des majeurs > 25 ans).
\begin{itemize}[label={}]
\item a) $\frac{1}{6}$ \quad b) $\frac{1}{4}$ \quad c) $\frac{1}{12}$ \quad d) $\frac{1}{3}$ \caseqcm
\end{itemize}

\item \textbf{Question :} Un débit de $36 \text{ m}^3\cdot\text{h}^{-1}$ correspond en $\text{L}\cdot\text{s}^{-1}$ à :
\begin{itemize}[label={}]
\item a) 3600 \quad b) 10 \quad c) 100 \quad d) 1000 \caseqcm
\end{itemize}

\item \textbf{Question :} Résoudre graphiquement $f(x) \le 5$ pour $f(x) = -x^2 + 10$.
\begin{center}
\begin{tikzpicture}[scale=0.6]
    \begin{axis}[axis lines=middle, xmin=-4.2, xmax=4.2, ymin=-2, ymax=11, 
        xtick={-2.23, 2.23}, xticklabels={$-\sqrt{5}$, $\sqrt{5}$}, ytick={5}, ylabel={$y$}, xlabel={$x$}]
        \addplot[blue, thick, domain=-3.5:3.5, samples=100] {-x^2 + 10};
        \addplot[red, dashed] coordinates {(-4,5) (4,5)};
    \end{axis}
\end{tikzpicture}
\end{center}
\begin{itemize}[label={}]
\item a) $[-\sqrt{5} ; \sqrt{5}]$ \quad b) $]-\infty ; -\sqrt{5}] \cup [\sqrt{5} ; +\infty[$ \quad c) $x \ge \sqrt{5}$ \caseqcm
\end{itemize}

\item \textbf{Question :} Soit la série statistique suivante :
\begin{center}
\renewcommand{\arraystretch}{1.5}
\begin{tabular}{|l|c|c|c|}
\hline
\textbf{Note} & 10 & 7 & $X$ \\ \hline
\textbf{Coefficient} & 1 & 2 & 2 \\ \hline
\end{tabular}
\end{center}
Que doit valoir $X$ pour que la moyenne soit égale à 12 ?
\begin{itemize}[label={}]
\item a) 17 \quad b) 18 \quad c) 19 \quad d) 20 \caseqcm
\end{itemize}

\item \textbf{Question :} Baisse de 30\% puis de 20\%. Baisse globale :
\begin{itemize}[label={}]
\item a) 10\% \quad b) 56\% \quad c) 44\% \quad d) 60\% \caseqcm
\end{itemize}

\item \textbf{Question :} Factoriser l'expression $x^2 - 16$.
\begin{itemize}[label={}]
\item a) $(x-8)^2$ \quad b) $(x-4)^2$ \quad c) $(x-4)(x+4)$ \quad d) $x(x-16)$ \caseqcm
\end{itemize}

\item \textbf{Question :} Quel est le coefficient multiplicateur d'une hausse de 0,5\% ?
\begin{itemize}[label={}]
\item a) 1,5 \quad b) 1,05 \quad c) 1,005 \quad d) 0,5 \caseqcm
\end{itemize}

\item \textbf{Question :} Résoudre l'équation $3x + 4 = 0$.
\begin{itemize}[label={}]
\item a) $x = 4/3$ \quad b) $x = -4/3$ \quad c) $x = -3/4$ \quad d) $x = 3/4$ \caseqcm
\end{itemize}

\item \textbf{Question :} L'image de -1 par la fonction $g(x) = 2x^2 - 3x + 1$ est :
\begin{itemize}[label={}]
\item a) 0 \quad b) 6 \quad c) -4 \quad d) 2 \caseqcm
\end{itemize}
\end{enumerate}

\newpage

\section*{Partie 2 : Exercices (14 points)}

\subsection*{EXERCICE 1 : Fonctions et Suites ($\approx$ 6 pts)}
$g(x) = \frac{3+x}{2x}$ et $f(x) = x-2$.

\textbf{1°) Montrer que $g(x)=f(x)$ revient à résoudre $\frac{-2x^2+5x+3}{2x} = 0$.} \reponseDouble

\textbf{2°) Donner les coordonnées des points d'intersection entre $C_g$ et $C_f$.} \reponse

\subsection*{EXERCICE 2 : Le Piston ($\approx$ 2 pts)}
$h(t) = 0,05 \cos(13t)$ et $v(t) = -0,65 \sin(13t)$.
\reponseDouble

\subsection*{EXERCICE 3 : Géométrie ($\approx$ 6 pts)}
$EFGH$ rectangle ($EH=2, EF=3$), $M$ milieu de $[FG]$, $\vec{HK} = \frac{1}{3}\vec{HG}$.

\begin{center}
\begin{tikzpicture}[scale=1.8]
    % Points
    \coordinate (E) at (0,0); \coordinate (F) at (3,0); 
    \coordinate (G) at (3,2); \coordinate (H) at (0,2);
    \coordinate (M) at (3,1); \coordinate (K) at (1,2);
    
    % Projection orthogonale de K sur (EM)
    \coordinate (L) at ($(E)!(K)!(M)$);
    
    % Dessin
    \draw[thick] (E) -- (F) -- (G) -- (H) -- cycle;
    \draw[blue, thick] (E) -- (M);
    \draw[red, thick] (E) -- (K);
    \draw[dashed] (K) -- (L);
    
    % Angle droit précis au point L
    \draw[thick] ($(L)!0.1!(E)$) -- ($(L)!0.1!(E)+(L)!0.1!(K)-(L)$) -- ($(L)!0.1!(K)$);
    
    % Labels
    \node[below left] at (E) {E}; \node[below right] at (F) {F};
    \node[above right] at (G) {G}; \node[above left] at (H) {H};
    \node[right] at (M) {M}; \node[above] at (K) {K};
    \node[below right] at (L) {L};
    
    % Cotes
    \draw[<->] (0,-0.3) -- (3,-0.3) node[midway, below] {3};
    \draw[<->] (-0.3,0) -- (-0.3,2) node[midway, left] {2};
\end{tikzpicture}
\end{center}

\textbf{1°) Montrer que $\vec{EK} \cdot \vec{EM} = 5$.} \reponseDouble

\textbf{2°) Déterminer la longueur $EL$ et une mesure de l'angle $\widehat{KEM}$.} \reponseDouble

\end{document}